\reportfigure{\begin{tikzpicture}
\begin{axis}[width=.95\linewidth,height=4.7cm,xmin=-340,xmax=340,
xtick={-300,-150,0,150,300},xlabel={Distance from fire axis (cm)},
ylabel={Temperature rise (K)},ymajorgrids,legend columns=3,
legend style={at={(.5,1.02)},anchor=south}]
\addplot[ink,only marks,mark=*,mark size=1.1pt] table[col sep=comma,x=figure_r_cm,y=exp_rise_K]{../../05_Calculations/op1_47_plume_diag_stations.csv};\addlegendentry{Peng et al.}
\addplot[teal,thick] table[col sep=comma,x expr=\thisrow{model_x_cm}-330,y=gas_rise_K]{../../05_Calculations/op1_47_plume_diag_stations.csv};\addlegendentry{FDS gas}
\addplot[rust,thick,dashed] table[col sep=comma,x expr=\thisrow{model_x_cm}-330,y=fds_rise_K]{../../05_Calculations/op1_47_plume_diag_stations.csv};\addlegendentry{FDS thermocouple}
\end{axis}\end{tikzpicture}}{Figure F1. OP-1 at 47.2 kW. Gas and 1 mm thermocouple temperatures were sampled together 20 mm below the ceiling on the adiabatic coarse grid. The 90--120~s time-weighted rises are compared with Peng et al.~\cite{peng}.}
